% filename: chap-experimental-execution-FIXED.tex
\chapter{Experimental Execution and Implementation}
\label{chap:experimental-execution}
\setlength{\emergencystretch}{1.5em}   % prevents right-margin overflow in this chapter

Chapter~\ref{chap:methodology} defined the experimental design: the
simulator architecture, policy set, and evaluation protocol.  The present
chapter documents how that design was executed in practice.  It records the
engineering decisions made during implementation, summarises the training
workload carried out across all sites and seeds, and explains how the final
ablation set evolved from the original experimental plan.  The objective is
to provide transparency on the experimental execution while avoiding
duplication of the methodological definitions already presented in
Chapter~\ref{chap:methodology}.


% ============================================================
\section{Implementation Overview}
\label{sec:impl_overview}
% ============================================================

The implementation was built in layers, each producing runnable and testable
output before the next layer began.  This strategy ensured that no layer
depended on untested components from a lower layer, and that partial results
were available at every stage.

\begin{table}[htb]
\centering
\small
\caption{Implementation execution sequence.}
\label{tab:impl_sequence}
\begin{tabular}{lll}
\toprule
\textbf{Layer} & \textbf{Component} & \textbf{Deliverable} \\
\midrule
1 & Data pipeline   & ITU/Zindi site CSVs loaded, validated, split \\
2 & Simulator       & \texttt{TelecomEnv} passing unit tests on all transitions \\
3 & Baselines       & B0, B1 rule-based policies producing first cost numbers \\
4 & RL agent        & RLInv training on single site converging \\
5 & Track B         & Hierarchical $(s,S)$ + PPO dispatch training \\
6 & Ablations       & A5, A6, A7 variants trained and evaluated \\
7 & Evaluation      & Full 1{,}800-rollout evaluation campaign complete \\
\bottomrule
\end{tabular}
\end{table}

The data pipeline and simulator (layers~1--2) were implemented and validated
before any policy training began.  Because all policies share the same
environment, correctness of the simulator was verified first to avoid
propagating modelling errors into downstream policy comparisons.
Rule-based baselines (layer~3) were implemented next because they require
no training and immediately produce interpretable cost numbers that serve as
sanity checks for the environment dynamics.

% FIX page 95: split the long path-string sentence across lines using \allowbreak
The repository structure follows the layer ordering:
\texttt{env/},
\texttt{baselines/},
\texttt{models/},
\texttt{train/},
\texttt{eval/},
with all hyperparameters in a single
\texttt{config/\allowbreak hparams.yaml}
and no magic numbers in training scripts.


% ============================================================
\section{Simulator Implementation Decisions}
\label{sec:impl_decisions}
% ============================================================

Several design decisions were made during simulator implementation that
diverge from or refine the formal model of Chapter~\ref{chap:problem}.
Each decision was made for a specific practical reason and is documented
here for reproducibility and examiner transparency.

\paragraph{Geometric rather than lognormal lead-time model.}
The formal model specified a lognormal lead-time distribution.  During
implementation, the geometric distribution was adopted instead.  The
geometric model preserves the Markov property: the probability of delivery
at each step is constant regardless of how long the order has been in
transit.  This means the pending-order state can be represented as a
two-scalar tuple (flag, quantity) rather than a countdown vector, keeping
the observation space low-dimensional.  The mean lead times are matched
between the two parameterisations (24\,h normal, 48\,h delayed), so the
operational stress is equivalent.

\paragraph{Discrete order actions.}
The formal model allowed a continuous order quantity.  The implementation
uses three discrete order levels: zero, $0.30\,C^{\Tank}$, and
$0.60\,C^{\Tank}$.  Discrete actions simplify the policy gradient
computation and match the physical reality of tanker deliveries, which
arrive in fixed load sizes rather than arbitrary quantities.  The reduction
from continuous to discrete was confirmed not to restrict meaningful
ordering behaviour: all three levels are exercised by the trained policies.

% FIX page 96: \allowbreak inside long monospace tokens prevents right-margin overflow
\paragraph{Action masking implementation.}
\texttt{Maskable\-PPO} from \texttt{sb3\-contrib} was used rather than
implementing a custom masking layer.  The library applies the mask by
setting infeasible action logits to $-\infty$ before the softmax, making
the probability of any infeasible action exactly zero by construction.
This required integrating a mask-computation function into the
environment's \texttt{step()} call and passing it through the
\texttt{Subproc\-Vec\-Env} wrapper --- a non-trivial integration verified
by checking that no illegal actions appeared in 10{,}000 random rollouts.

\paragraph{Tank capacity derivation.}
The ITU dataset does not provide absolute diesel tank capacity.  Tank
capacity was set to $C^{\Tank} = 72\,\bar{D}$ kWh-equivalent, where
$\bar{D}$ is the site mean hourly load, giving approximately three days
of full-load backup.  This convention was applied consistently across all
ten sites and produces physically plausible values.

\paragraph{Observation normalisation synchronisation.}
\texttt{Vec\-Normalize} accumulates running mean and variance statistics
during training.  If these statistics are not explicitly frozen and
synchronised before evaluation, the normalisation applied to evaluation
observations differs from that seen during training, causing a systematic
input-scale mismatch.  A custom evaluation callback was implemented to
copy the training \texttt{Vec\-Normalize} state to the evaluation
environment before each validation pass.  Without this fix, early
evaluation results showed anomalously high variance across seeds that
disappeared once the synchronisation was corrected.
\sloppy
\paragraph{Parallel environment count.}
The training configuration targeted eight parallel workers using
\texttt{Subproc\-Vec\-Env}, following the architecture plan.  On
the Windows development machine, the OS process-handle limit caused
intermittent failures at eight workers.  The parallel environment count
was reduced to four workers for all training runs.  This doubles
wall-clock training time per run but does not affect the statistical
properties of the PPO gradient estimates, as the effective batch size is
preserved by proportionally increasing $n_{\text{steps}}$.
\fussy
% ============================================================
\section{Training Execution}
\label{sec:training_execution}
% ============================================================

\subsection{Workload summary}

Table~\ref{tab:training_workload} summarises the complete training
workload executed across all policies, sites, and seeds.

% FIX page 97: \small + narrower Note column prevents table overflow
\begin{table}[htb]
\centering
\small
\setlength{\tabcolsep}{5pt}
\caption{Total training workload.}
\label{tab:training_workload}
\begin{tabular}{lrp{5.5cm}}
\toprule
\textbf{Item} & \textbf{Count} & \textbf{Note} \\
\midrule
Sites                    & 10  & All ITU/Zindi sites \\
Seeds per policy         & 3   & Seeds 42, 123, 777 \\
RL policies trained      & 4   & RLInv, TrackB, A5, A6 \\
Non-RL policies          & 3   & B0, B1 (no training); A7 (PPO, no masking) \\
Timesteps per RL run     & 400{,}000 & Per site per seed \\
Total RL training runs   & 120 & 4 policies $\times$ 10 sites $\times$ 3 seeds \\
A7 targeted rerun        & 6   & Sites 5, 10 $\times$ 3 seeds \\
Total training runs      & 126 & Including A7 rerun \\
Approx.\ wall-clock time & $\approx$\,2\,h & Per run, 4-core CPU \\
\bottomrule
\end{tabular}
\end{table}

\subsection{Convergence check}

A convergence check was conducted on three representative sites (1, 5,
and~7) using seed~42, extending training to 560{,}000 steps.  Sites~1
and~7 showed no change in evaluation reward beyond 400{,}000 steps.
Site~5 exhibited high reward variance across independent runs regardless
of training length, consistent with its classification as a
grid-constrained site where episode outcomes are sensitive to the
stochastic lead-time draws.  The 400{,}000-step checkpoint was retained
for all sites and all policies.  Extended training was not pursued
further because the variance on site~5 was attributed to environment
stochasticity rather than insufficient policy training.

\subsection{Training stability}

RLInv and TrackB training was stable across all sites and seeds.  A5
(no inventory observation) showed higher early-episode variance on Hard
sites, consistent with the expectation that removing inventory
information degrades ordering decisions.  A7 (no action masking)
produced unstable training curves on sites~5 and~10, with occasional
reward spikes followed by policy collapse.  For consistency, the final
model checkpoint (rather than a best-episode checkpoint) was used for
all policies.


% ============================================================
\section{Ablation Set Evolution}
\label{sec:ablation_evolution}
% ============================================================

The original experimental plan specified eight ablations labelled A1--A8
following the architecture plan of Section~2 of the project
specification.  The final executed ablation set differs from this plan.
Table~\ref{tab:ablation_evolution} documents every change and the reason
for it.

\begin{table}[htb]
\centering
\small
\caption{Evolution from planned ablation set (A1--A8) to executed set.}
\label{tab:ablation_evolution}
\begin{tabular}{llp{7.0cm}}
\toprule
\textbf{Original label} & \textbf{Final label} & \textbf{Decision and reason} \\
\midrule
A1 (Full SC-PPO)  & RLInv   & Renamed to reflect actual implementation
                               (MaskablePPO flat MLP, not modular SC-PPO) \\
A8 (Hierarchical) & TrackB  & Renamed; calibrated $(s,S)$ + PPO dispatch \\
A5 (No Inv state) & A5      & Retained unchanged \\
A7 (Vanilla PPO)  & A7      & Retained; evaluated on constrained sites only
                               (sites~5 and~10) where masking effect is measurable \\
B0, B1            & B0, B1  & Rule-based baselines added (not in original A1--A8
                               plan); essential anchors for cost comparison \\
\midrule
A2 (No masking alone) & Dropped & A7 already tests the no-masking condition;
                                  A2 was redundant given MaskablePPO design \\
A3 (No Lagrangian)    & Dropped & Lagrangian update was not implemented in the
                                  final system (masking alone drove violations
                                  to zero); A3 has no executable meaning \\
A4 (No recovery)      & Dropped & Recovery policy was not implemented
                                  (masking handles the cases it was designed
                                  for); A4 has no executable meaning \\
A6 (Fixed ordering)   & A6 added & Not in original plan; introduced after
                                   the TrackB comparison was found to be
                                   potentially confounded, providing a
                                   cleaner H1 test by holding dispatch
                                   architecture constant \\
\midrule
\textbf{Final set} & \multicolumn{2}{l}{B0, B1, A5, A6, A7, TrackB, RLInv} \\
\bottomrule
\end{tabular}
\end{table}

The most important change is the addition of A6.  When TrackB results
became available, the comparison between RLInv and TrackB was potentially
confounded because TrackB uses a different dispatch action space in
addition to a different ordering mechanism.  A6 was therefore introduced
to isolate the ordering contribution specifically by holding the dispatch
architecture constant while varying only the ordering mechanism.


% ============================================================
\section{Evaluation Execution}
\label{sec:eval_execution}
% ============================================================

\subsection{Evaluation rollout count}

The evaluation campaign comprised the following rollout workload:

\[
30 \text{ episodes} \times 10 \text{ sites} \times 3 \text{ seeds}
\times 2 \text{ scenarios} = 1{,}800 \text{ rollouts for RLInv and TrackB}
\]

Baselines and ablations were evaluated under the normal scenario only,
giving an additional $30 \times 10 \times 3 = 900$ rollouts per policy.
The total evaluation dataset contains 11{,}160 episode records across all
policies, sites, seeds, and scenarios.

\subsection{Initial inventory randomisation}

Training episodes initialised inventory at a fixed $0.6\,C^{\Tank}$ to
reduce early-episode variance and allow the policy to observe a range of
ordering decisions before inventory reaches critical levels.  Evaluation
episodes randomise initial inventory uniformly over
$[0.3,\,0.9]\,C^{\Tank}$, testing the policy across a realistic range
of starting conditions.  This asymmetry between training and evaluation
initial conditions is a deliberate design choice: a policy that trains
only at $0.6\,C^{\Tank}$ must generalise to lower and higher starting
inventories at evaluation time, which is the realistic deployment
scenario.

\subsection{Evaluation on constrained sites}

The A7 ablation (vanilla PPO without masking) was evaluated only on
sites~5 and~10, the most reliability-constrained sites in the dataset.
Initial trials on less constrained sites did not expose the effect of
removing action masking, because battery autonomy and outage severity
were not sufficient to force frequent infeasible-action situations.
Sites~5 and~10 were therefore selected for the targeted rerun because
their battery autonomy (3.8\,h and 12.1\,h respectively) and outage
frequency (48.3\% both) make constraint violations much more likely
without masking.


% ============================================================
\section{Metrics and Statistical Framework}
\label{sec:metrics}
% ============================================================

\subsection{Operational metrics}

Table~\ref{tab:metrics} defines the primary metrics computed from each
evaluation episode.

\begin{table}[htb]
\centering
\small
\caption{Evaluation metrics and their operational interpretation.}
\label{tab:metrics}
\begin{tabular}{lp{4.2cm}p{4.8cm}}
\toprule
\textbf{Metric} & \textbf{Definition} & \textbf{Interpretation} \\
\midrule
EENS (kWh)      & Expected Energy Not Served per episode: $\sum_t U_t$
                & Primary reliability metric.  Zero means
                  all load served; higher values indicate
                  more unserved load. \\[4pt]
Diesel consumption (kWh) & Total DG energy dispatched per episode
                & Proxy for fuel cost and carbon
                  emissions. \\[4pt]
Stockout events & Number of hours with $\Inv_t = 0$
                  and $G_t = 0$
                & Counts hard reliability failures;
                  directly corresponds to service
                  interruptions. \\[4pt]
Mean inventory (\%) & Mean $\Inv_t / C^{\Tank}$ per episode
                & Indicates ordering behaviour;
                  very low values suggest under-ordering. \\[4pt]
Orders placed   & Count of non-zero order actions per episode
                & Characterises ordering frequency. \\
\bottomrule
\end{tabular}
\end{table}

EENS is the primary metric throughout the results analysis.  It
directly measures the consequence of policy failure --- unserved telecom
load --- and is sensitive to both dispatch and inventory decisions.
Diesel consumption and stockout events are reported as secondary metrics
to distinguish policies that achieve similar EENS through different
mechanisms (e.g.\ high diesel use vs.\ efficient battery management).

\subsection{Statistical unit and testing}

The statistical unit of analysis is the seed-level mean: for each
policy, site group, and scenario, the 30 episode values within a seed
are averaged to produce one observation per seed, giving $n = 3$
observations per condition.  Using the episode-level mean as the unit
($n = 30$) would be incorrect because the 30 episodes within a seed
share the same trained weights --- they are not independent observations
of a random policy.  The seed-level mean correctly treats each
independently trained policy as one observation.

Statistical comparisons use Welch's $t$-test (unequal variances) with
Cohen's $d$ reported alongside $p$-values to provide an effect-size
estimate in addition to significance.  Given $n = 3$ per group, the
$t$-test has low power against small effects; results are therefore
interpreted in conjunction with the observed percentage difference and
effect size rather than relying on $p$-values alone.  All tests are
two-sided with no correction for multiple comparisons, as the three
hypotheses (H1, H2, H3) are pre-registered and tested independently.